
\def \Subject {مبانی و نظریه}
\def \Session { یک}
% \setcounter{chapter}{\Session-1}
\renewcommand{\chaptername}{سوال}
\chapter{\Subject}
\section{بخش اول}
\subsection{ پرسش ۱}

\subsubsection{الف}

\begin{enumerate}
	\item 
	\textbf{استقلال:}
	این شرط بیان می‌کند که باید مقدار پیش‌بینی‌شده نسبت به متغیر تصادفی زیرگروه‌ها مستقل باشد. یعنی اگر در نظر بگیریم 
	\lr{A}
	متغیر تصادفی زیرگروه و 
	\lr{R}
	متغیر تصادفی مربوط به مقدار پیش‌بینی‌شده باشد، داریم:
	\begin{latin}
		~\\
		$
		A \perp R \Rightarrow P(R | A) = P(R) \Rightarrow P(R | A = a) = P(R | A = b) = P(R)
		$
	\end{latin}
	یکی از محدودیت‌های این معیار این است که حتی اگر زیرگروه 
	\lr{A}
	تاثیر بدون سوگیری روی نتیجه داشته باشد، ما آن تاثیر را نیز در نظر نمی‌گیریم. مثلا جنسیت می‌تواند روی سابقه‌کار و سابقه‌کار روی درصد پذیرش رزومه تاثیر داشته باشد، ولی اگر  معیار هدف استقلال باشد،‌این اطلاعات به کلی نادیده گرفته می‌شود. 
	\textbf{جداسازی:}
	این شرط بیان می‌کند که مقدار پیش‌بینی‌شده باید با توجه به مقدار واقعی برچسب، نسبت به زیرگروه‌ها مستقل باشد. به عبارت دیگر، اگر
	\lr{A}
	متغیر تصادفی زیرگروه،
	\lr{Y}
	برچسب واقعی و
	\lr{R}
	مقدار پیش‌بینی‌شده باشد، خواهیم داشت:
	
	\begin{latin}
		~\\
		$
		R \perp A \mid Y
		\Rightarrow
		P(R \mid A,Y)=P(R\mid Y)
		$
		
		$
		P(R\mid A=a,Y=y)=P(R\mid A=b,Y=y)
		$
	\end{latin}
	
	این معیار به نوعی تا حدی مشکل معیار قبل را سعی می‌کند حل کند تا اطلاعات تاثیرگذار روی خروجی از دست نروند. برای همین مقدار برچسب را در شرط قرار می‌دهد.
	
	این معیار تضمین می‌کند که برای افرادی که برچسب واقعی یکسانی دارند، توزیع خروجی مدل در تمامی زیرگروه‌ها یکسان باشد. بنابراین مدل نباید برای افراد با وضعیت واقعی یکسان، صرفاً به دلیل تعلق به یک زیرگروه خاص، رفتار متفاوتی داشته باشد. 
	
	یک شیوه نگاه به این معیار این است که نرخ خطا در زیرگروه‌های مختلف باید یکی باشد. اگر مدلی که احتمال وقوع جرم را پیش‌بینی می‌کند، روی افراد سیاه‌پوست نرخ خطای بیشتری دارد می‌توان گفت مدل دچار سوگیری است. 
	
	این معیار اجازه می‌دهد توزیع کلی پیش‌بینی‌ها میان زیرگروه‌ها متفاوت باشد، به شرط آنکه این تفاوت ناشی از اختلاف در برچسب‌های واقعی باشد.
	
	\vspace{0.5cm}
	
	\textbf{بسندگی:}
	این شرط بیان می‌کند که برچسب واقعی باید با توجه به مقدار پیش‌بینی‌شده، نسبت به زیرگروه‌ها مستقل باشد. اگر
	\lr{A}
	متغیر تصادفی زیرگروه،
	\lr{Y}
	برچسب واقعی و
	\lr{R}
	مقدار پیش‌بینی‌شده باشد، خواهیم داشت:
	
	\begin{latin}
		~\\
		$
		Y \perp A \mid R
		\Rightarrow
		P(Y\mid A,R)=P(Y\mid R)
		$
		
		$
		P(Y\mid A=a,R=r)=P(Y\mid A=b,R=r)
		$
	\end{latin}
	
	این معیار بیان می‌کند که مقدار پیش‌بینی‌شده باید در تمامی زیرگروه‌ها تفسیر یکسانی داشته باشد. به عبارت دیگر، اگر دو فرد از دو زیرگروه مختلف امتیاز یا احتمال پیش‌بینی یکسانی دریافت کنند، احتمال واقعی تعلق آن‌ها به کلاس مثبت نیز باید یکسان باشد. یعنی پیش‌بینی مدل اطلاعاتی راجع به زیرگروه ندهد.
\end{enumerate}

\subsubsection{ب}

برای هر یک از سه معیار آماری، یک معیار انصاف متناظر برای تصمیم‌گیری دودویی به صورت زیر تعریف می‌شود:

\begin{enumerate}
	\item \textbf{استقلال: برابری جمعیتی}
	
	این معیار بیان می‌کند که نرخ پذیرش باید در تمام زیرگروه‌ها یکسان باشد. بنابراین
	
	\begin{latin}
		\[
		P(\hat{Y}=1 \mid A=a)=P(\hat{Y}=1 \mid A=b)
		\]
	\end{latin}
	
	به عبارت دیگر، احتمال دریافت تصمیم مثبت نباید به عضویت در زیرگروه وابسته باشد.
	
	\item \textbf{جداسازی: برابری شانس‌ها}
	
	این معیار بیان می‌کند که برای افراد با برچسب واقعی یکسان، خروجی مدل نباید به زیرگروه وابسته باشد. بنابراین
	
	\begin{latin}
		\[
		P(\hat{Y}=1\mid Y=y,A=a)
		=
		P(\hat{Y}=1\mid Y=y,A=b),
		\qquad y\in\{0,1\}
		\]
	\end{latin}
	
	که معادل برابر بودن نرخ مثبت واقعی 
	\LTRfootnote{True positive rate}
	و نرخ مثبت کاذب
	\LTRfootnote{False positive rate}
	در تمام زیرگروه‌ها است.
	
	\item \textbf{بسندگی: برابری پیش‌بینی}
	
	این معیار بیان می‌کند که اگر مدل تصمیم مثبت بگیرد، احتمال درست بودن این تصمیم باید در همه زیرگروه‌ها یکسان باشد. بنابراین
	
	\begin{latin}
		\[
		P(Y=1\mid \hat{Y}=1,A=a)
		=
		P(Y=1\mid \hat{Y}=1,A=b)
		\]
	\end{latin}
	
	در حالت پیش‌بینی احتمالاتی، این معیار معادل کالیبراسیون است.
\end{enumerate}


\subsubsection{ج}

\begin{enumerate}
	\item \textbf{استقلال}
	
	مثلا می‌توان نرخ قبولی رزومه را بررسی کرد.  هدف این معیار آن است که نرخ دعوت به مصاحبه میان زنان و مردان یکسان باشد تا جنسیت به طور مستقیم بر تصمیم اثر نگذارد.
	
	\item \textbf{جداسازی}
	
	
	مثلا در یک سیستم پیش‌بینی تکرار جرم مهم است که نرخ خطا برای تمام گروه‌های نژادی یکسان باشد تا یک گروه بیش از گروه دیگر به اشتباه پرخطر تشخیص داده نشود.
	
	\item \textbf{بسندگی}
	
	مثلا در یک سیستم که صلاحیت افراد برای دریافت وام را بررسی می‌کند، اگر دو متقاضی از دو گروه مختلف امتیاز اعتباری ۸۰ درصد دریافت کنند، بانک باید بتواند این امتیاز را برای هر دو به یک معنا تفسیر کند؛ در غیر این صورت امتیاز مدل قابل اعتماد نیست.
\end{enumerate}

\subsubsection{د}

فرض کنید نرخ پایه در دو زیرگروه متفاوت باشد، یعنی

\begin{latin}
	\[
	P(Y=1\mid A=0)\neq P(Y=1\mid A=1).
	\]
\end{latin}

همچنین فرض کنید مدل غیر بدیهی باشد؛ یعنی پیش‌بینی آن نه کاملاً صحیح باشد و نه مستقل از برچسب واقعی.
اگر مقدار 
$
P(Y | A) 
$
را با استفاده از قانون احتمال کل برای متغیر 
\lr{R}
بنویسیم داریم:
\begin{latin}
	~\\
	$
	P(Y | A ) = \sum\limits_r P(Y | A, R) P(R | A) \xrightarrow{P(R | A ) = P(R)}  
	P(Y | A) = \sum\limits_r P(Y | A, R) P(R) \xrightarrow{P(Y | A, R) = P(Y | R)} 
	P(Y | A) =  \sum\limits_r P(Y | R) P(R)
	$
\end{latin}
می‌توان دید سمت راست تساوی هیچ اثری از 
\lr{A}
وجود ندارد پس 
$P(Y | A) = P(Y)$
ولی این با فرض اولیه که در نظر گرفته بودیم 
$P(Y | A = a) \neq P(Y | A = b)$
در تناقض است. پس در حالتی که نرخ‌های پایه باهم تفاوت دارند نمی‌توان این دو معیار را باهم داشت. 

\subsection{پرسش ۲}
\subsubsection{الف}
\begin{latin}
	~\\
	$
	PPV = P(Y = 1 | \hat{Y} = 1) = \cfrac{TP}{TP + FP} \Rightarrow 1 - PPV = \cfrac{FP}{TP + FP} 
	\\ \Rightarrow
	\frac{1 - PPV}{PPV} = \frac{FP}{TP} \xrightarrow{1 - FNR =\frac{TP}{TP + FN}} \frac{1 - PPV}{PPV} . (1 - FNR)  = 
	\cfrac{FP}{TP} . \cfrac{TP}{TP + FN} 
	= 
	\cfrac{FP}{TP + FN}
	\xrightarrow{p = \cfrac{TP + FN}{TP + TN + FP + FN}}
	\cfrac{p}{1-p}. \cfrac{1 - PPV}{PPV} . (1 - FNR) = \cfrac{TP + FN}{TN + FP}.  \cfrac{FP}{TP + FN}
	= \cfrac{FP}{TN + FP} = FPR
	% 
	$
\end{latin}
\subsubsection{ب}
با استفاده از الف و فرض ثابت بودن 
\lr{PPV}
داریم:
\begin{latin}
	~\\
	$
	\cfrac{FPR}{1 - FNR} = \cfrac{p}{1 - p} * constant
	$
\end{latin}
حال اگر مقدار 
\lr{FPR}
و
\lr{FNR}
یکی باشد پس سمت چپ تساوی برای هر دو گروه یکی می‌شود پس مقدار
\lr{p}
نیز باید برای هر دو گروه یکی شود که در تناقض با فرض است.

حالاتی که می‌توان این مقادیر یکی شود وقتی است که در اتحاد ضرب یا تقسیم در صفر داشته باشیم. یعنی مقدار
\lr{PPV}
برابر صفر یا یک شود و یا مقدار 
\lr{FNR}
برابر یک شود.

\subsubsection{ج}
(برای شناخت مسئله گفته‌شده در سوال از هوش مصنوعی و همچنین مقاله 
\cite{compass_debate}
استفاده شد.
)
شرکت سازنده معیار 
\lr{sufficiency }
را هدف قرار داد. یعنی اگر یک فرد سفیدپوست و یک فرد سیاه‌پوست ریسک یکسانی بگیرند احتمال تکرار جرم هردوی آن‌ها یکسان است. ولی 
\lr{ProPublica}
نشان داد که نرخ خطا بین زیرگروه‌ها متفاوت است. 

شباهت این موضوع با سوال در این است که اگر معیار 
\lr{sufficiency }
درست باشد پس انگار مقدار 
\lr{PPV}
در بین زیرگروه‌ها یکی است. پس مشابه بخش «ب» از آنجایی که نرخ پایه در بین زیرگروه‌ها متفاوت است، پس نرخ خطا نیز متفاوت است.

پس یعنی اگر نرخ‌های پایه متفاوت باشد نمی‌توان همزمان هم معیار 
\lr{sufficiency}
و
\lr{separation}
را باهم داشت.

\subsubsection{د}
اگر نرخ‌های پایه برابر باشد، تنها حالتی که همزمان با 
\lr{sufficiency}
بتوان تساوی بخش «الف» را برقرار کرد وقتی است که مقدار
\lr{FPR}
و
\lr{FNR}
هردو صفر باشد. پس اگر مدل کامل باشد می‌توان. ولی این حالت تقریبا هیچوقت رخ نمی‌دهد. زیرا فارغ از حجم داده خود مسئله معمولا یک نویز غیرقابل پیش‌بینی نیز دارد. در بهترین حالت می‌توان گفت هرچقدر مدل کامل‌تر باشد، از دید 
\lr{separation}
مدل به سمت \lr{fairness} حرکت می‌کند.
\subsection{پرسش سوم. انصاف علی و خلاف‌واقع}
\subsubsection{الف}
اگر U متغیرهای 
\lr{Exogenous}
و
V 
متغیرهای 
\lr{Endogenous}
باشد و A متغیر حساس باشد و مدل علی ساختاری زیر را داشته باشیم:

\begin{latin}
	~\\
	$
	M = (U, V, F)
	$
\end{latin}
آنگاه ایده خلاف واقع این است که برای یک نمونه ورودی با استفاده از عملگر do مقدار A آن را تغییر داد و باقی مقادیر U را ثابت نگه داشت و دید خروجی مدل تغییر می‌کند یا خیر. یعنی:

\begin{latin}
	~\\
	$
	P(\hat{Y}(U) = y | X = x, do(A = a)) =P(\hat{Y}(U) = y | X = x, do (A = a'))
	$
\end{latin}
\subsubsection{ب}
در اینجا متغیر
X
هم می‌تواند تبیین‌کننده باشد و هم پروکسی. اگر فقط مقدار و معنای X را در نظر بگیریم، این یک متغیر تبیین‌کننده است زیرا سال‌های تجربه روی امتیاز استخدام تاثیر دارد و سوگیری خاصی ندارد.

ولی ممکن است از روی سال‌های تجربه بتوان جنسیت را تخمین زد و این موضوع روی امتیاز استخدام اثر بگذارد. در اینصورت پروکسی رخ داده. 

مسیر تبیین کننده 
$X -> Y$
و مسیر پروکسی 
$A -> X -> Y$
است. باید از $X$ اطلاعات صرفا مربوط به سابقه را استخراج کرد و نه جنسیت.

اگر ما هیچ توضیحی از A X و Y نداشته باشیم نمی‌توان فهمید که روابط بین این متغیرهای تصادفی ناشی از سوگیری است و یا رابطه واقعی. برای مثال اگر در گراف سوال مقدار A را برابر سخت‌کوشی در نظر بگیریم همین گراف بدست می‌آید ولی این گراف دیگر سوگیری ندارد چون سخت‌کوشی یک متغیر تبیین‌کننده است. برای همین تشخیص اینکه چه متغیرهایی دچار سوگیری مدل می‌شوند و این که مدل چه معیاری از عدالت را رعایت کند تصمیمی است که از بیرون و براساس هنجار گرفته می‌شود و نه از روی داده.

\subsubsection{ج}
زیرا همانطور که در بخش قبل ذکر شد اینکه کدام متغیر تاثیر درست دارد و کدام متغیر سوگیری، براساس هنجار است و نه داده. حتی اگر توضیح هر متغیر موجود باشد، ممکن است یک نفر جنسیت را متغیر دارای سوگیری نداند و در نظر بگیرد جنسیت برای استخدام یک متغیر درست و بدون سوگیری است. مثلا برای موقعیت شغلی مربی مهدکودک معمولا خانم‌ها قبول می‌شوند.

از این رو اگر یک نفر جنسیت را متغیر دارای اثر بداند یال 
$A -> \hat{Y}$
را در نظر می‌گیرد. ولی اگر جنسیت را سوگیری در نظر بگیریم این یال نباید در گراف علیت باشد.

\section{بخش دوم. پرسش‌های مفهومی و کاربردی}
\subsection{پرسش ۴: مسئله جانشین/برچسب در سلامت}
\subsubsection{الف}
(برای فهمیدن انواع بایاسی که در سوال ذکر شد از هوش مصنوعی استفاده شد.)

\begin{itemize}
	\item 
	\textbf{تاریخی:}
	وجود دارد. گروه‌هایی وجود دارد که در گذشته دسترسی نابرابر به خدمات درمانی داشتند و یا هزینه درمان برای آن‌ها متفاوت‌تر است.
	\textbf{بازنمایی و تجمیع و ارزیابی}
	در اینجا خیلی مسئله نیست ولی در کل احتمالا وجود دارد.
	\textbf{اندازه‌گیری برچسب:}
	این مسئله اصلی است. برای اینکه تخمین قابل سنجش‌پذیرتر باشد مدل هزینه درمان را پیش‌بینی می‌کند و فرض کرده‌است بین هزینه و نیاز رابطه واضحی است. ولی هزینه به‌طور کامل نیاز را اندازه‌گیری نمی‌کند.
	\textbf{استقرار:}
	وجود دارد. مدل برای هدفی متفاوت با آنچه که استفاده می‌شد ساخته شد.
\end{itemize}

\subsubsection{ب}
خیر. زیرا در اینجا بیشتر از اینکه مدل روی یک ویژگی دچار سوگیری شده باشد، هدف مدل دارای سوگیری است. یعنی فارغ از هر ویژگی‌ای که مدل با آن آموزش ببیند، این فرض که نیاز از هزینه بدست می‌آید دارای سوگیری است. (هرچند که حتی اگر سوگیری در داده یا گروه بود نیز احتمالا حذف ویژگی مشکل را حل نمی‌کرد.)

\subsubsection{ج}
اولا باید نسبت به متغیر هدف اصلی یعنی نیاز پزشکی سنجید. زیرا سوگیری درون هدف جایگزین یعنی هزینه است. پس اگر نسبت به هزینه حساب کنیم، احتمالا مشکلی پیدا نکنیم.

در درجه دوم مقاله معیار 
\lr{separation}
را حساب می‌کند. یعنی بررسی می‌کند در بیماران با سطح بیماری یکسان کدام گروه نرخ خطای بیشتری دارد.  (هرچند احتمالا اگر با روشی متفاوت معیارهای دیگر حساب می‌شد نیز سوگیری در آن معلوم بود.)

\subsubsection{د}
انگیزه اصلی این که بجای نیاز هزینه پیش‌بینی شد، این است که هزینه قابل سنجش‌تر است. یک راه این است که مستقیم خود نیاز را ارزیابی کرد و از معیارهای عدالت برای سنجش سوگیری مدل استفاده کرد. ولی حتی اگر بخواهیم مدل قابل سنجش‌تری داشته باشیم، اولا هزینه خود تابع نیاز است و نه بالعکس و فقط اگر هزینه خام (بدون بیمه و نرمال شده با عوامل محیطی مثل محل و دکتر) را در نظر بگیریم می‌توان بررسی کرد که تا چه اندازه هزینه از روی نیاز می‌تواند تخمین زده شود. 

یعنی حتی اگر هزینه را همچنان به عنوان هدف جایگزین در نظر بگیریم باید عوامل دیگری که روی نیاز اثر می‌گذارند را حذف کرد.
\subsection{پرسش۵. انتخاب معیار، تحت قید عدم امکان (اعطای وام)}
\subsubsection{الف}
\begin{itemize}
	\item 
	\textbf{بانک: }
	معیار مهم برای بانک مقدار نرخ مثبت کاذب است. یعنی مدل یکسری وام را تایید کند ولی افراد بازپرداخت نداشته باشند. این مورد باعث ضرر بانک می‌شود. عامل کمتر مهم نرخ منفی کاذب است. این باعث نمی‌شود که بانک ضرر کند ولی باعث می‌شود که سود کمتری کند زیرا می‌توانست با دادن این وام‌ها سود بیشتری بکند. 
	
	\item 
	\textbf{مشتریان:}
	معیار مهم برای مشترین  مقدار نرخ منفی کاذب است. زیرا باعث می‌شود تا وام نگیرند در حالی که توانایی بازپرداخت وام را داشته‌اند. 
	
	\item 
	\textbf{ناظران:}
	باتوجه به اینکه این ناظران قاعده ۸۰ درصد را مدنظر قرار داده‌اند، یعنی به نرخ پذیرش در گروه‌های مختلف حساسند. هرچقدر گروه‌های مختلف نرخ قبولی متفاوت‌تری داشته باشند، مسئله بزرگ‌تر است. پس یعنی هم منفی کاذب و هم مثبت کاذب برای این افراد مهم است. (درواقع اصلا این که پیش‌بینی مدل چقدر با واقعیت نزدیک است برای این گروه مهم نیست و فقط برابری گروه‌ها مسئله است.)
\end{itemize}
\subsubsection{ب}
باتوجه به اینکه در سوال قاعده ۸۰ درصد ذکر شد پس برابری جمعیتی معیار نزدیک‌تری به نظر می‌رسد. این مورد اطمینان حاصل می‌کند که خروجی مدل از گروه افراد مستقل است. پس نرخ قبولی گروه‌های مختلف تقریبا برابر است و قاعده ۸۰ درصد که بیان می‌کند کمترین گروه باید حداقل ۸۰ درصد بیشترین گروه باشد رعایت می‌شود.

البته این موضوع بدون هزینه نیست. اگر نرخ‌های پایه برابر باشد مسئله‌ای نیست ولی اگر برابر نباشد چون 
\lr{seperation}
نداریم پس یعنی نرخ خطای گروه‌های مختلف برابر نیست. یعنی ممکن است یک گروه مقدار منفی‌های کاذب بیشتری داشته باشد صرفا به این دلیل که نرخ قبولی آن با بقیه گروه‌ها یکی شود. پس یعنی گروه روی خروجی تاثیر می‌گذارد. ولی حداقل می‌توان گفت که مدل ترکیب جمعیتی را تغییر نمی‌دهد. 

\subsection{پرسش ۶. داوری مناقشه COMPASS}
\subsubsection{الف}
این مسئله به این صورت است که شرکت 
NorthPointe 
یک مدل ارائه داد که برای مجرمین احتمال تکرار جرم توسط آن‌ها را پیش‌بینی می‌کرد. این شرکت برای عدالت بیان کرد که اگر دو فرد با دو نژاد مختلف ریسک یکسانی بگیرند احتمال واقعی تکرار جرم آن‌ها مشابه است. یعنی معیار 
\lr{sufficiency}
را مبنای عدالت قرار داد. اما 
\lr{ProPublica}
روی نرخ خطا در بین نژادهای مختلف تمرکز کرد و نشان داد که نرخ خطا در بین افراد سفیدپوست و سیاه‌پوست یکی نیست و افراد سیاه‌پوست نرخ خطای بیشتری دارند.

یعنی شرکت می‌گفت که پیش‌بینی دارد بدون سوگیری سعی می‌کند مقدار واقعی را پیش‌بینی کند. ولی 
ProPublica
می‌گفت که اگر مقدار واقعی این‌که جرم تکرار شده یا نه را بدانیم، و مدل را با آن تست کنیم، افراد سیاه‌پوست نرخ خطای بیشتری خواهند داشت. 

\subsubsection{ب}
همانطور که در بخش‌های قبل توضیح داده شد، شرکت عملا گفته مقدار 
PPV
در بین نژادهای مختلف یکی است. ولی در همان سوال قبل نیز دیدیم که اگر نرخ پایه متفاوت باشد و PPV یکسان باشد، مقدار نرخ خطا نیز متفاوت می‌شود. پس یعنی امکان‌پذیر نیست که در هر دو معیار عدالت داشت. 

در حقیقت چون مدل معیار 
\lr{sufficiency}
را رعایت کرده مجبور است معیار 
\lr{seperation}
را نقض کند. هیچ کدام از طرفین دروغ نمی‌گفتند بلکه هرکدام معیار عدالت را چیز متفاوتی درنظر گرفته بودند. 
\subsubsection{ج}
به طور کلی در سیستم‌های قضائی ترجیح این است که بجای اینکه فردی بصورت اشتباه محاکمه شود، یک فرد گناهکار محاکمه نشود. یعنی اگر متغیر تصادفی Y گناهکار بودن را نشان بدهد،‌اولویت بر کاهش 
FP 
ها است و سپس کاهش FN ها.

در این مسئله اگر هدف این باشد که مدل میزان ریسکی که پیش‌بینی می‌کند واقعا فارغ از نژاد احتمال تکرار جرم را نشان بدهد، معیار \lr{seperability} مهم است. و اگر هدف برابری کامل باشد \lr{seperation} من به نظرم در اینجا علیرغم بحثی که راجع به \lr{false positive} شد، این که ریسک فارغ از گروه مفهوم یکسانی را بدهد و بتوان به آن اعتماد کرد مهم‌تر است. زیرا تصمیم نهایی را قاضی و یا هیئت منصفه می‌گیرند و این ریسک به نوعی نشان می‌دهد که تا چه میزان دقت نیاز است. پس یعنی حتی اگر مدل ریسکی را به اشتباه بالا تخمین بزند و مثبت کاذب ایجاد کند، این باعث می‌شود تا سیستم قضائی بیشتر دقت کند ولی اگر مدرک مستدلی یافت نشد فرد تبرئه می‌شود.  (البته باید دقت داشت که توزیع جامعه به هم نخورد.)

ولی اگر قرار است که خروجی مدل مستقیم منجر به محکومیت و یا تبرئه فرد شود، \lr{seperation} احتمالا معیار عادلانه‌تری است. همچنین نابرابری‌ای که در سوال‌های قبل بدست آمد در حالت مطلق است. یعنی نمی‌توان تماما هر دو معیار را داشت ولی می‌توان ترکیبی از هر دو معیار را مدنظر گرفت.

\subsubsection{د}
خیر. زیرا سوگیری‌های موجود در داده حاصل از نتایج دنیای واقعی هستند. پس یعنی فارغ از مدل این سوگیری‌ها در دنیا وجود دارند. خواه بخاطر محیط باشد یا سطح تمکن افراد و یا ذهنیت افراد. پس اگر مدل را کامل حذف کنیم مسئله کامل حل نمی‌شود. همچنین ذکر شد که از آنجایی که مدل تصمیم نهایی را نمی‌گیرد، این ملاحظه که نباید مثبت کاذب داشته باشیم همچنان رعایت می‌شود. یعنی به نوعی مدل مقدار منفی کاذب را کاهش می‌دهد و بجای این‌که مثبت کاذب زیاد شود فقط سیستم قضائی هزینه‌های بیشتری متحمل می‌شود (زیرا تحقیقات جامع‌تری انجام داده است.)

\subsection{پرسش ۷. فراتر از انصاف حاشیه‌ای - زیرگروه‌ها و زمان‌ها}
\subsubsection{الف}
مثلا فرض کنید یک مدل تایید وام داریم که خروجی آن در جدول
\ref{tab:bg}
آمده است. اگر فقط به توزیع‌های حاشیه‌ای نگاه کنیم، نرخ تایید مدل برای هر نژادی برابر ۵۰ درصد و برای هر جنسیتی نیز برابر ۵۰ درصد است که برابر با نرخ تایید کل است و به ظاهر سوگیری‌ای نداریم ولی اگر دقیق‌تر نگاه کنیم می‌بینیم در زیرگروه‌های توام مدل دچار سوگیری شده است. 

یک راه می‌تواند این باشد که همه زیرگروه‌ها را امتحان می‌تواند این باشد که مدل را به ازای زیرگروه‌های مختلف امتحان کرد. ولی تعداد زیرگروه‌ها برابر با 
$3^d - 1$
است زیرا هر ویژگی یا مثبت است یا منفی و یا بی تاثیر. پس یعنی حتی با ابعاد کم هم امکان بررسی همه زیرگروه‌ها خیلی ممکن نیست، چه برسد به اینکه ابعاد زیاد شود.

        \begin{table}[htbp] \label{tab:bg}
	\centering
	\caption{ نتایج طبقه‌بندی مدل دارای سوگیری}
	\begin{tabular}{|c|c|c|c|}
		\hline
		\textbf{نژاد} & \textbf{جنسیت} & \textbf{تعداد} & \textbf{خروجی مدل} \\
		سیاه پوست & مرد & ۱۰۰ & ۱ \\
				سفید پوست & مرد  & ۱۰۰ & ۰ \\
						سیاه پوست & زن  & ۱۰۰ & ۰ \\
								سفید پوست & زن & ۱۰۰  & ۱ \\
								\hline
	\end{tabular}
\end{table}

\subsubsection{ب}
اگر یک مدل تایید وام داشته باشیم، به دلیل 
\lr{fairness}
ما وام افراد یک گروه که امتیاز پایین گرفته‌اند را به دلیل اینکه نرخ قبولی آن‌ها نیز برابر باقی شود تایید می‌کنیم. منتها این مدل درواقع احتمال بازپرداخت وام را حساب می‌کرده. پس افرادی که امتیاز پایین گرفته‌اند احتمال بازپرداخت وام توسط آن‌ها پایین است و با پرداخت نکردن وام، دچار آسیب بیشتری می‌شوند.  

























\section{بخش سوم. پیاده‌سازی و تحلیل تجربی}
(در این سوال، برای تبدیل جداول به جدول 
\lr{latex}
و بعضا نوشتن یک سری تابع کمکی برای رسم نمودار و جدول از هوش مصنوعی استفاده شد.
)


\subsection{بررسی اولیه‌ی داده‌ها}
نرخ پایه برای هریک از نژادها در جدول 
\ref{tab:race}
آمده است. 
        \begin{table}[htbp] \label{tab:race}
        	\centering
        	\begin{latin}


\begin{tabular}{|lr|}
	\toprule
	race & two\_year\_recid \\

	African-American & 0.523150 \\
	Caucasian & 0.390870 \\
	\bottomrule

\end{tabular}
	        	\end{latin}
\end{table}




توزیع متغیر \lr{decile\_score} برای گروه‌های نژادی مختلف در شکل \ref{fig:score_distribution} آمده است و نشان می‌دهد که توزیع امتیاز خطر در میان گروه‌های مختلف یکسان نیست و برخی گروه‌ها به‌طور متوسط امتیازهای بالاتری دریافت کرده‌اند. چنین اختلافی نشان‌دهنده‌ی تفاوت در توزیع داده‌ها میان گروه‌های حساس است که می‌تواند منجر به ایجاد سوگیری در مدل‌های پیش‌بینی شود.

\begin{figure}[H]
	\centering
	\includegraphics[width=0.75\linewidth]{images/compas_decile_score_distribution_by_race.png}
	\caption{توزیع امتیاز خطر COMPAS بر حسب نژاد}
	\label{fig:score_distribution}
\end{figure}

به‌منظور بررسی میزان اطلاعات نژادی موجود در ویژگی‌ها، یک مدل کمکی برای پیش‌بینی نژاد از روی سایر ویژگی‌ها آموزش داده شد. این مدل توانست با دقت حدود $0.67$ مقدار نژاد را پیش‌بینی کند. 

مهم‌ترین ویژگی‌های مؤثر در این پیش‌بینی در شکل \ref{fig:proxy_features} نمایش داده شده‌اند. مشاهده می‌شود که حتی در صورت حذف مستقیم ویژگی نژاد، برخی ویژگی‌ها همچنان نقش متغیرهای جانشین (\lr{Proxy Features}) را ایفا می‌کنند. این موضوع یکی از دلایل اصلی باقی ماندن سوگیری پس از حذف ویژگی حساس است. 

پس یعنی نمی‌توان گفت که صرفا با حذف ویژگی نژاد مسئله سوگیری حل می‌شود و با راه‌های نسبتا پیچیده‌تری باید این مسئله را حل کرد.

\begin{figure}[H]
	\centering
	\includegraphics[width=0.8\linewidth]{images/top_proxy_features_for_race_prediction.png}
	\caption{مهم‌ترین ویژگی‌های جانشین برای پیش‌بینی نژاد}
	\label{fig:proxy_features}
\end{figure}

\begin{table}[H]
	\centering
	\caption{گزارش عملکرد مدل برای پیش‌بینی ویژگی نژاد از روی سایر ویژگی‌ها}
	\label{tab:classification_report_race}
	\begin{latin}


	\begin{tabular}{|lcccc|}
		\toprule
		\textbf{Group} & \textbf{Precision} & \textbf{Recall} & \textbf{F1-score} & \textbf{Support} \\
		\midrule
		African-American & 0.69 & 0.83 & 0.75 & 635 \\
		Caucasian        & 0.62 & 0.44 & 0.52 & 421 \\
		\midrule
		Accuracy          & \multicolumn{3}{c}{0.67} & 1056 \\
		Macro Avg         & 0.66 & 0.63 & 0.63 & 1056 \\
		Weighted Avg      & 0.66 & 0.67 & 0.66 & 1056 \\
		\bottomrule
	\end{tabular}
	\end{latin}
\end{table}

\subsection{مدل پایه}

در ادامه یک مدل Logistic Regression بر روی داده‌های آموزش آموزش داده شد. منحنی ROC مدل در شکل \ref{fig:roc} نمایش داده شده است. مقدار بالای سطح زیر منحنی (AUC) نشان می‌دهد که مدل از قدرت تفکیک مناسبی برخوردار است و می‌تواند افراد دارای احتمال تکرار جرم را از سایر افراد با دقت قابل قبول متمایز کند.

\begin{figure}[H]
	\centering
	\includegraphics[width=0.7\linewidth]{images/roc_curve.png}
	\caption{منحنی ROC مدل پایه}
	\label{fig:roc}
\end{figure}

نتایج مربوط به 
\lr{fairness}
در جدول \ref{tab:fairness_metrics} آمده است. می‌توان دید که به غیر از 
\lr{demographic parity}
مدل در سایر زمینه‌ها عملکرد بسیار خوبی دارد. مثلا هر دو نژاد شانس برابری دارند. به این معنی که نرخ خطا برابر است. پس در ادامه بررسی‌ها هدف بهبود 
\lr{demographic parity}
است.

\begin{table}[H]
	\centering
	\caption{اختلاف معیارهای انصاف میان گروه‌های نژادی}
	\label{tab:fairness_metrics}
	
	\begin{latin}
		\begin{tabular}{|lc|}
			\hline
			\textbf{Metric} & \textbf{Between-group Gap} \\
			\hline
			
			Demographic Parity & 0.141 \\
			
			Equal Opportunity & 0.000 \\
			
			False Positive Rate Parity Difference (FPR Gap) & 0.042 \\
			
			PPV Difference & 0.015 \\
			\hline
		\end{tabular}
	\end{latin}
	
\end{table}



\subsection{تحلیل عملکرد مدل بر حسب نژاد}

برای بررسی دقیق‌تر رفتار مدل، منحنی ROC برای هر گروه نژادی به‌صورت جداگانه رسم شد. نتیجه در شکل \ref{fig:roc_race} ارائه شده است.

\begin{figure}[H]
	\centering
	\includegraphics[width=0.75\linewidth]{images/roc_curves_by_race.png}
	\caption{منحنی ROC به تفکیک گروه‌های نژادی}
	\label{fig:roc_race}
\end{figure}

همان‌گونه که مشاهده می‌شود، منحنی مربوط به افراد سفیدپوست  در تمامی نقاط بالاتر از منحنی افراد سیاه‌پوست قرار دارد و در مقادیر نسبتاً کوچک FPR به مقدار تقریباً یک برای TPR می‌رسد. در مقابل، منحنی ROC مربوط به افراد سیاه‌پوست فاصله‌ی بسیار کمتری با خط قطری دارد و در نتیجه توان تفکیک مدل برای این گروه ضعیف‌تر است.

این اختلاف نشان می‌دهد که مدل اطلاعات موجود در داده‌های مربوط به گروه سفیدپوست را بسیار بهتر از گروه سیاه‌پوست یاد گرفته است. از آنجا که 
\lr{Equalized Odds }
مستلزم برابری همزمان TPR و FPR میان گروه‌ها است، فاصله‌ی زیاد میان منحنی‌های ROC بیانگر آن است که دستیابی به این معیار تنها از طریق تغییر آستانه‌ی تصمیم دشوار خواهد بود و احتمالاً مستلزم کاهش قابل توجهی در دقت مدل است.

\subsection{روش‌های کاهش سوگیری}

به‌منظور مقایسه‌ی سه دسته‌ی اصلی روش‌های کاهش سوگیری، سه رویکرد مختلف مورد استفاده قرار گرفت.

در مرحله‌ی پیش‌پردازش از الگوریتم \lr{Disparate Impact Remover} استفاده شد. این روش با تغییر توزیع برخی ویژگی‌ها تلاش می‌کند وابستگی میان ویژگی‌های ورودی و ویژگی حساس را کاهش دهد، بدون آنکه برچسب‌های داده تغییر کنند.

در مرحله‌ی درون‌پردازش از الگوریتم \lr{Exponentiated Gradient} همراه با قید \lr{Demographic Parity} استفاده شد. این روش هنگام آموزش مدل، قید انصاف را مستقیماً وارد مسئله‌ی بهینه‌سازی می‌کند و میان بیشینه‌سازی دقت و رعایت معیار انصاف نوعی مصالحه ایجاد می‌نماید.

در مرحله‌ی پس‌پردازش  بهینه‌سازی آستانه چک شد. منتها همانطور که در شکل 
\ref{fig:roc_race}
آمده است، 
\lr{ROC curve}
ها از هم واگرا هستند، پس با هر استانه‌ای باز نمی‌توان به میزان خوبی از 
\lr{fairness}
رسید، برای همین این موضوع بیش از این بررسی نشد. (زیرا سایر روش‌ها دقت بهتری می‌دهند. )

\begin{table}[H]
	\centering
	\caption{مقادیر سنجه‌های انصاف برای حذف کننده تاثیر نامتناسب}
	\label{tab:fairness_dir}
	
	\begin{latin}
		\begin{tabular}{|lc|}
			\hline
			\textbf{Metric} & \textbf{Between-group Gap} \\
			\hline

			Demographic Parity & 0.123 \\

			Equal Opportunity & 0.028 \\

			False Positive Rate Parity Difference (FPR Gap) & 0.035 \\

			PPV Difference & 0.012 \\
			\hline
		\end{tabular}
	\end{latin}
	
\end{table}

\begin{table}[H]
	\centering
	\caption{مقادیر سنجه‌های انصاف برای روش ‌گرادیان نمایی   }
	\label{tab:fairness_expgrad}
	
	\begin{latin}
		\begin{tabular}{|lc|}
			\hline
			\textbf{Metric} & \textbf{Between-group Gap} \\
			\hline
			
			Demographic Parity & 0.028 \\
			
			Equal Opportunity & 0.093 \\
			
			False Positive Rate Parity Difference (FPR Gap) & 0.049 \\
			
			PPV Difference & 0.098 \\
			\hline
		\end{tabular}
	\end{latin}
	
\end{table}

\begin{table}[H]
	\centering
	\caption{دقت مدل‌های مختلف}
	\label{tab:model_accuracy}
	
	\begin{latin}
		\begin{tabular}{|lc|}
			\hline
			\textbf{Model} & \textbf{Accuracy} \\
			\hline
			
			Baseline Logistic Regression & 0.97 \\
			
			Pre-processing (Disparate Impact Remover) & 0.96 \\
			
			In-processing (Exponentiated Gradient) & 0.88 \\
			\hline
		\end{tabular}
	\end{latin}
	
\end{table}

\subsection{بده‌بستان میان دقت و انصاف}

به‌منظور بررسی اثر تغییر آستانه، مقدار آستانه‌ی تصمیم در ۹ مقدار مختلف بین ۰٫۱ تا ۱ تغییر داده شد. برای هر مقدار، Accuracy و اختلاف 
\lr{Demographic Parity }
و 
\lr{Equal Opportunity}
محاسبه گردید. نتیجه در شکل \ref{fig:tradeoff} نشان داده شده است.

\begin{figure}[H]
	\centering
	\includegraphics[width=0.75\linewidth]{images/accuracy_vs_fairness_tradeoff.png}
	\caption{بده‌بستان میان دقت و انصاف}
	\label{fig:tradeoff}
\end{figure}

نتایج نشان می‌دهد که با افزایش Accuracy، اختلاف Demographic Parity نیز افزایش پیدا می‌کند. به بیان دیگر، هرچه مدل از نظر دقت بهتر عمل می‌کند، نرخ پذیرش میان گروه‌های مختلف نیز از یکدیگر فاصله‌ی بیشتری می‌گیرد. این موضوع بیانگر وجود یک بده‌بستان میان بیشینه‌سازی عملکرد مدل و رعایت Demographic Parity است.

در مقابل، هنگام بررسی معیار Equal Opportunity روندی تقریباً معکوس مشاهده شد؛ به‌گونه‌ای که بهبود این معیار معمولاً با کاهش Accuracy همراه بود. این نتیجه با مباحث نظری مطرح‌شده در بخش دوم همخوانی دارد و نشان می‌دهد که معیارهای مختلف انصاف لزوماً در یک جهت تغییر نمی‌کنند و بهبود یکی از آن‌ها ممکن است موجب بدتر شدن معیار دیگر شود.

پس یعنی ما هرچقدر سعی در بهبود معیار
\lr{Demographic Parity}
بکنیم، مقدار
\lr{Equal Opportunity}
کاهش می‌یابد. این مسئله در بین مدل‌های مختلف نیز مشهود است. برای مثال در مدل پایه همانطور که در جدول 
\ref{tab:fairness_metrics}
می‌توان دید، مقدار
\lr{Equal Opporunity}
ایده‌آل است. ولی در مدل گرادیان نمایی، مقدار
\lr{Demographic Parity}
خیلی بهتر شده است ولی این به قیمت کمی بدتر شدن 
\lr{Equal Opportunity }
تمام شده است.

\subsection{بررسی اتحاد Chouldechova}

برای هر یک از آستانه‌های تصمیم، مقادیر PPV، FPR و FNR برای هر گروه نژادی محاسبه و اتحاد Chouldechova بررسی شد. اختلاف میان مقدار واقعی PPV و مقدار محاسبه‌شده از رابطه‌ی نظری در تمامی آزمایش‌ها  در بیشترین حالت حدود
$5.56e-17
$
ود. که این اختلاف خیلی ناچیز نیز می‌تواند به علت نویز ذاتی داده و همچنین اختلاف تخمین احتمال از روی داده باشد.


\subsection{تحلیل تقاطعی}

عملکرد مدل گرادیان نمایی برای گروه‌های تقاطعی حاصل از ترکیب نژاد و جنسیت نیز بررسی شد. نتایج مقادیر انصاف در جدول 
\ref{tab:intersectional_fairness}
آمده است. این مقادیر به ازای گروه‌های نژاد، جنسیت و نژاد-جنسیت حساب شده‌اند. همانطور که می‌توان دید، مدل هم در بخش نژاد و هم در بخش جنسیت از لحاظ انصاف خوب است. ولی وقتی به ترکیب دو ویژگی نگاه می‌کنیم مدل عملکرد متوسطی دارد.

\begin{table}[H]
	\centering
	\caption{سنجه‌های انصاف برای گروه‌های حساس مختلف در تحلیل تقاطعی}
	\label{tab:intersectional_fairness}
	
	\begin{latin}
		\begin{tabular}{|l|l|c|}
			\hline
			\textbf{Sensitive Attribute} & \textbf{Metric} & \textbf{Between-group Gap} \\
			\hline
			
			\multirow{4}{*}{Race}
			& Demographic Parity & 0.046 \\
			& Equal Opportunity & 0.075 \\
			& False Positive Rate Parity Difference (FPR Gap) & 0.031 \\
			& PPV Difference & 0.083 \\
			\hline
			
			\multirow{4}{*}{Sex}
			& Demographic Parity & 0.031 \\
			& Equal Opportunity & 0.061 \\
			& False Positive Rate Parity Difference (FPR Gap) & 0.044 \\
			& PPV Difference & 0.094 \\
			\hline
			
			\multirow{4}{*}{Race--Sex}
			& Demographic Parity & 0.066 \\
			& Equal Opportunity & 0.109 \\
			& False Positive Rate Parity Difference (FPR Gap) & 0.089 \\
			& PPV Difference & 0.174 \\
			\hline
		\end{tabular}
	\end{latin}
	
\end{table}




